\cotefile{144}
\scanpage{122}
On s'intéresse à des objets du type suivant :
\[
\left\{\begin{array}{l}
\text{a) une variété abélienne à isogénie près } A\text{ sur }\overline{\mathbb{Q}},\text{ de dim }g\\[4pt]
\text{b) une structure de }K\text{-module sur }A\\[4pt]
\text{c) pour chaque }\ell,\text{ un isomorphisme de }K\text{-modules}
\end{array}\right.
\]
\diagram{\begin{tikzcd}[column sep=large,row sep=small]
& V_\ell(A) \arrow[d,equal] \\
K\otimes\mathbb{Q}_\ell \arrow[r,leftrightarrow,"\unsure{t_\ell|\iota_\ell}"] & T_\ell(A)\otimes\mathbb{Q}_\ell
\end{tikzcd}}
et on veut déterminer comment $\operatorname{Gal}(\overline{\mathbb{Q}}/\mathbb{Q})$ permute les classes d'isomorphie de tels objets. On n'y arrive pas tout à fait.

Sur $\overline{\mathbb{Q}}$ (ie après extension des scalaires à $\overline{\mathbb{Q}}$), la $\mathbb{Q}$-algèbre $K$ devient un produit de copies de $\overline{\mathbb{Q}}$ :
\[
K\otimes\overline{\mathbb{Q}}\xrightarrow{\sim}\overline{\mathbb{Q}}^{\operatorname{Hom}(K,\overline{\mathbb{Q}})}
\]
Notons $K^*$ le groupe algébrique ``groupe multiplicatif de la $\mathbb{Q}$-algèbre $K$ (c'est $R_{K/\mathbb{Q}}(\mathbb{G}_m)$). $\operatorname{Lie}(A)$ est une représentation de ce groupe -- sur $\overline{\mathbb{Q}}$.
De même pour $\bigwedge^g\operatorname{Lie}(A)$, qui définit, sur $\overline{\mathbb{Q}}$, un caractère
\[
\det\operatorname{Lie}:K^*\longrightarrow\mathbb{G}_m\qquad(\text{défini sur }\overline{\mathbb{Q}}).
\]
Le \uline{corps dual} $K'\subset\overline{\mathbb{Q}}$ est le corps de définition de ce caractère, et on n'arrive à décrire que l'action de $\operatorname{Gal}(\overline{\mathbb{Q}}/K')$.

J'aurais aussi bien pu définir $K'$ comme le plus petit sous-corps de $\overline{\mathbb{Q}}$ tel que le $K$-$\overline{\mathbb{Q}}$-bimodule $\operatorname{Lie} A$ provienne, par extension des scalaires, d'un $K$-$K'$-bimodule $L_0$ (unique à isomorphisme non unique près) : $\operatorname{Lie} A\sim L_0\otimes_{K'}\overline{\mathbb{Q}}$

Ce bimodule détermine un morphisme de groupes algébriques $/\mathbb{Q}$ :
\[
\delta:K'{}^*\longrightarrow K^*:\quad k'\longmapsto\det_K(\text{multiplication par }k'),
\]
qui induit à son tour
\[
\delta:\text{ classes d'idèles de }K'{}^*\longrightarrow\text{ classes d'idèles de }K
\]
\margin{122}
